\section{Related Work}
\label{sec:related}
\paragraph{Minimal assumptions behind a verdict.}
Bastani et al.~\cite{bastani-oopsla15} compute min-cuts in a source-sink flow graph, each cut a set of statements assumed unreachable, and ask a developer to confirm them with tests; Zhu et al.~\cite{zhu-aplas13} abduce the library specifications needed to verify a source-sink property; Dillig et al.~\cite{dillig-pldi12} and Elrond~\cite{elrond} abduce proof obligations or specifications.
Our receipt is their cut with prices; the human becomes optional, a calibrated prior carries the rest with a certified error, and when a person does check facts, receipts tell which few facts unlock the most verdicts (RQ9).
\paragraph{Library specification inference.}
Merlin~\cite{merlin} infers specifications with probabilities but thresholds them into facts; Seldon~\cite{seldon}, StubDroid~\cite{stubdroid}, Atlas~\cite{atlas} and Bastani et al.~\cite{bastani-popl15} produce specifications without a verdict-level error; LLM-based extraction~\cite{iris,semtaint} takes the model's answers as fact.
We use StubDroid as an independent key and the model-as-fact approach as a baseline.
\paragraph{Probabilistic inference.}
Praline~\cite{praline} computes sound probability bounds for Datalog outputs under partly known input correlations; on a reachability fragment with only marginals known, its optimum equals the min-cut (Proposition~\ref{prop:tight}), which we use because it is a max-flow.
\paragraph{Neuro-symbolic analysis.}
NeuroLog~\cite{neurolog} feeds model-extracted facts to Datalog to find bugs; we use the same kind of stack to bound missed flows rather than to find them.
\paragraph{Risk control.}
The calibration is Learn-then-Test~\cite{lt}; conformal selection~\cite{conformalselection} is the verdict-level alternative; it controls the false-discovery rate among selected verdicts, which here leaves wrong clearances in most apps.
\paragraph{Benchmarks.}
TaintBench~\cite{taintbench} documents real malware leaks with their paths; we use it to ask how many of them depend on library holes and whether a threshold certified on benign apps transfers.
\paragraph{Soundiness.}
The soundiness manifesto~\cite{soundiness} asks analyses to list what they leave out; receipts make one such list checkable per verdict and attach a risk to it.

\section{Conclusion}
\label{sec:conclusion}
An unsound static analysis can be $(\alpha,\varepsilon)$-soundy: it can say on which library assumptions each clean verdict rests, and clear verdicts with an error rate over programs that is certified, not hoped for.
The certificate holds for every pricer we tried; a good pricer buys clearance, a few well-chosen checks of receipt facts buy more, and on real malware the threshold certified on benign apps transferred.
The limits are as clear: our analysis is a weaker tool than FlowDroid, the key is only partly independent of language models, and cheap models price too confidently to help.
